\section{Conclusion}\label{sec:conclusion}

We have established a mathematical framework connecting the algebraic
structure of error propagation to convergence behavior in multi-step
mathematical proofs.  The Lyapunov exponent ordering theorem
(Theorem~\ref{thm:lyapunov}) shows that abelian error operators
accumulate errors strictly more slowly than non-abelian ones, with the
gap controlled by the commutator Frobenius norm at order $\Theta(\eps^2)$.
The spectral gap comparison (Theorem~\ref{thm:spectral_gap}) reveals
that non-abelian groups can mix up to four times faster on matched Cayley
graphs, creating a fundamental tension between error minimization and
exploration efficiency.  The Diaconis bound analysis
(Theorem~\ref{thm:diaconis}) explains why abelian groups yield more
informative convergence bounds.  Together, these results provide a
principled basis for proof strategy design: commuting substeps for error
control, non-commuting operations for search efficiency, and the
commutator norm as a computable diagnostic.

These results bridge Lie group error theory \cite{barrau2014,wang2008,li2023},
quantum error correction \cite{gottesman1996,arvind2002,zhao2023},
and spectral theory of random walks \cite{diaconis1988,furstenberg1963}
into a unified framework applicable to mathematical proof systems.
